\section{Evaluation Design and Results}

All computational methods receive the same discovery and held-out evaluation task
manifests, generator-model and request-randomness policy, independent Oracle and
functional contract, candidate budget, family-valid intervention protocol,
randomization, outcome encoding, semantic-cluster uncertainty, and
multiplicity policy. Candidate universes, selector scores and ranks, native
candidates, mappings, expected directions, and realization policies are frozen
before held-out outcomes exist.

Observational selection uses the discovery split, and randomized effect
evaluation uses held-out tasks. The primary confirmatory analysis is hypothesis- and
generator-model-specific semantic-task-clustered assigned-arm ITT on oracle-evaluable
secure-code yield. Secure-and-functional joint success is the key practical
secondary outcome. Every
committed assignment remains in its assigned-arm denominator;
\texttt{target\_changed}, semantic compliance, and non-target drift are
diagnostics rather than filters.

Throughout the tables, \texttt{--} marks a result slot awaiting
provenance-verified frozen-run or study backfill; it denotes neither zero, not
applicable, nor observed missingness.

\subsection{Study Data}

The task inventory draws from SecurityEval, LLMSecEval, SALLM, CWEval,
CyberSecEval, and SeCodePLT, with BaxBench used for backend replication
\cite{siddiq2022securityeval,tony2023llmseceval,siddiq2024sallm,
peng2025cweval,bhatt2023cyberseceval,nie2025secodeplt,vero2025baxbench}.
The frozen task manifest identifies the exact included records and preserves
dataset provenance and semantic-cluster-disjoint study splits.

\subsection{RQ1: Held-Out Intervention Selection}

\textbf{RQ1. Under a fixed top-$K$ budget, how effectively do different
selectors identify intervention hypotheses with supported effects on held-out
tasks?}

The primary selector-only comparison gives our approach, Pearson correlation,
elastic-net logistic regression, and LLM-based explanation the same shared
candidate universe, information budget, and $K$ slots. Pearson correlation
ranks each binary prompt factor by its absolute correlation with the frozen binary
security outcome. Elastic-net logistic regression models all prompt factors
jointly, selecting regularization strength by semantic-task-clustered
cross-validation within the discover split and ranking standardized coefficient
magnitudes. The LLM baseline explains and ranks the same candidates from the
same discovery evidence without access to held-out outcomes. Seeded random
ranking provides a chance reference. The unique union of the methods' top-$K$
skeletons crosses the common
bridge once and each resulting hypothesis is randomized once; duplicate
selector references reuse that evidence. Empty, bridge-failed, and
protocolization-failed slots remain zeroes in denominator $K$. The primary
endpoint is Confirmed Effect Yield@$K$, accompanied by simultaneous paired
selector-utility differences. Rankings remain conditional on the frozen
discover split; held-out outcomes neither rerank selectors nor replace failed
slots.

A separate native-system track evaluates Ji et al.'s causality-aided
prompt--code analysis and NLPerturbator in their native candidate spaces
\cite{ji2025causality,chen2026nlperturbator}. Their outputs are mapped to the
editable factor catalog before held-out randomized evaluation and follow
the fixed funnel

\[
K \rightarrow N_{\mathrm{native}} \rightarrow N_{\mathrm{mapped}}
\rightarrow N_{\mathrm{protocol}} \rightarrow N_{\mathrm{randomized}}
\rightarrow N_{\mathrm{confirmed}}.
\]

That track measures system compatibility rather than pure selector
superiority. Conditional mapping coverage and mapped yield at $K$ are bridge
diagnostics; native funnel results cannot replace the shared-universe primary
comparison. For every frozen method/slot hypothesis, the evidence report
retains target and operation, target-minus-no-op ITT, simultaneous status,
independent semantic-cluster and arm counts, and preregistered specificity
contrasts.

\begin{table}[H]
  \centering
  \small
  \caption{RQ1 shared-universe selector results. Panel A retains every frozen
  budget slot through effect evaluation. Panel B reports the corresponding
  yields; Confirmed Effect Yield@$K$ is primary. Our approach and the three
  baselines use identical candidate and information budgets.}
  \label{tab:rq1}
  \begin{tabular}{lrrrrrr}
    \toprule
    \multicolumn{7}{l}{\textit{Panel A: Frozen-slot funnel}} \\
    \midrule
    Selector & $K$ & Selected & Bridged & Protocol & Rand. & Conf. \\
    \midrule
    \model & -- & -- & -- & -- & -- & -- \\
    Pearson correlation & -- & -- & -- & -- & -- & -- \\
    Elastic-net logistic & -- & -- & -- & -- & -- & -- \\
    LLM-based explanation & -- & -- & -- & -- & -- & -- \\
    \bottomrule
  \end{tabular}

  \medskip
  \begin{tabular}{lrrrr}
    \toprule
    \multicolumn{5}{l}{\textit{Panel B: Derived yields}} \\
    \midrule
    Metric & \model & Pearson & Elastic net & LLM expl. \\
    \midrule
    Selected / $K$ & -- & -- & -- & -- \\
    Bridged / $K$ & -- & -- & -- & -- \\
    Protocol / $K$ & -- & -- & -- & -- \\
    Unique-union share & -- & -- & -- & -- \\
    Block-freeze coverage & -- & -- & -- & -- \\
    Randomized / $K$ & -- & -- & -- & -- \\
    Confirmed Effect Yield@$K$ & -- & -- & -- & -- \\
    Confirmed / Randomized & -- & -- & -- & -- \\
    \bottomrule
  \end{tabular}
\end{table}
\begin{table}[H]
  \centering
  \small
  \caption{RQ1 hypothesis and paired-selector evidence. The frozen RQ1
  manifest binds selector/slot references to one shared randomized hypothesis
  estimate. Paired utility intervals resample top-level semantic clusters and
  never rerun discovery.}
  \label{tab:rq1-evidence}
  \begin{tabular}{llllrr}
    \toprule
    Selector/slot & Hyp.\ ID & Frozen target & Op. & ITT RD & Interval \\
    \midrule
    \multicolumn{6}{c}{\emph{Rows are populated only from the
    provenance-verified frozen RQ1 manifest}} \\
    \bottomrule
  \end{tabular}

  \medskip
  \begin{tabular}{llrrl}
    \toprule
    Selector pair & Utility diff. & Simult.\ interval & Status & Shared slots \\
    \midrule
    \multicolumn{5}{c}{\emph{Continued fields for the same frozen rows}} \\
    \bottomrule
  \end{tabular}
\end{table}

\subsection{RQ2: Contribution of Context-Conditioned Representation}

\textbf{RQ2. How does context conditioning in the Prompt TSG representation
affect candidate validity, protocolization, and held-out selection utility?}

RQ1 already isolates selector differences on one shared universe. RQ2 instead
holds the selector, budget, bridge, arm policies, outcome construction, and
inference fixed while comparing a direct-actionable-feature universe with the
context-conditioned Prompt TSG universe. Candidate identities and eligible
populations may differ, so this is an end-to-end representation comparison,
not a second selector contest.

The comparison reports candidate coverage, context applicability,
invalid-operation rate, protocolization rate, randomized coverage, Confirmed
Effect Yield@$K$, and the distribution of hypothesis-specific effects. The
context-conditioned representation is valuable when it converts broad factor
names into operations that are applicable, task-preserving, and evaluable for
a stated security context. No selected hypotheses are pooled into an
unapproved cross-hypothesis ATE.

\begin{table}[H]
  \centering
  \scriptsize
  \caption{RQ2 compares direct features with context-conditioned Prompt TSG
  templates end to end. All non-representation coordinates are frozen; the
  comparison is not interpreted as pure selector superiority.}
  \label{tab:rq2}
  \resizebox{\linewidth}{!}{%
  \begin{tabular}{lrrrrrr}
    \toprule
    Representation & Candidates & Context appl. & Invalid op. & Protocol & Rand. & Effect Yield@$K$ \\
    \midrule
    Direct features & -- & -- & -- & -- & -- & -- \\
    Prompt TSG context & -- & -- & -- & -- & -- & -- \\
    \bottomrule
  \end{tabular}}
\end{table}

\begin{table}[H]
  \centering
  \small
  \caption{RQ2 hypothesis-specific evidence retains each representation
  universe and context-conditioned population. Every row reports the
  target-minus-no-op ITT without cross-hypothesis pooling.}
  \label{tab:rq2-evidence}
  \begin{tabular}{llllrr}
    \toprule
    Universe & Hyp.\ ID & Context / target & Op. & ITT RD & Interval \\
    \midrule
    \multicolumn{6}{c}{\emph{Rows are populated only from the
    provenance-verified frozen RQ2 manifest}} \\
    \bottomrule
  \end{tabular}

  \medskip
  \begin{tabular}{llrrl}
    \toprule
    Hyp.\ ID & Effect status & Indep.\ clusters & Arm counts & Specificity \\
    \midrule
    \multicolumn{5}{c}{\emph{Continued fields for the same frozen rows}} \\
    \bottomrule
  \end{tabular}
\end{table}

\Needspace{8\baselineskip}
\subsection{RQ3: Security Intervention Effects}

\textbf{RQ3. Which frozen prompt interventions improve or degrade secure-code
yield, and what are their specificity and functionality profiles?}

RQ3 estimates context-conditioned, generator-model-specific, multi-realization policy
effects for the frozen security targets. The primary outcome is
oracle-evaluable secure-code yield; secure-and-functional joint success is the
key practical secondary outcome. Code validity, Oracle support, unknowns,
terminal no-code rates, and secure-yield bounds remain separate so that a
coverage or production shift cannot masquerade as safer generated code.

ADD and REMOVE remain separate candidate slots, hypotheses, protocols,
contrasts, multiplicity coordinates, and evidence. Observed baseline or no-op
insecurity is neither an eligibility condition nor an ITT denominator filter.
Target-versus-placebo and, when family-valid, target-versus-generic contrasts
test specificity. Bidirectional support is awarded only when separately
randomized ADD and REMOVE policies support opposite expected directions.

\begin{table}[H]
  \centering
  \small
  \caption{RQ3 reports hypothesis-specific target-versus-no-op assigned-arm
  effects on oracle-evaluable secure-code yield. Effect, specificity, and
  functionality statuses remain separate so that mixed evidence is visible.}
  \label{tab:rq3}
  \scriptsize
  \resizebox{\linewidth}{!}{%
  \begin{tabular}{lllrrlll}
    \toprule
    Intervention & Op. & Generator & ITT RD & Interval & Effect status & Specificity & Functionality \\
    \midrule
    \multicolumn{8}{c}{\emph{Rows are populated only from the
    provenance-verified frozen TargetSpec manifest}} \\
    \bottomrule
  \end{tabular}}
\end{table}

Target-change, semantic-compliance, non-target-drift, task improvement, harm,
and tie rates are descriptive diagnostics. They neither filter the assigned-arm
ITT denominator nor replace the primary ITT estimates.

\subsection{RQ4: Expert-Perceived Explanation Utility}

\textbf{RQ4. How do security experts rate and rank the perceived quality and
usefulness of intervention explanations produced by different methods?}

For \model, the intervention evidence record is rendered as an expert-facing
explanation containing the context, proposed edit, observational selection
trace, randomized effect and interval, specificity, functionality, and claim
scope. Each method's explanation is placed in the same anonymous wrapper and
presented under a balanced case--method assignment. Qualified experts rate clarity,
evidence sufficiency, credibility, actionability, and overall usefulness on
five-point ordinal scales. Overall usefulness is the preregistered primary
human-study outcome; the other dimensions are secondary.

The primary analysis uses an ordinal mixed-effects model with method as a fixed
effect and participant and case as random effects, followed by
multiplicity-adjusted comparisons. RQ4 concerns perceived utility only and
does not claim objective factor identification, repair accuracy, or
validation of the computational causal results.

\begin{table}[H]
  \centering
  \small
  \caption{RQ4 expert-rating results. Entries report ordinal summaries and
  model-based contrasts from the preregistered expert study.}
  \label{tab:rq4}
  \begin{tabular}{lrrrrr}
    \toprule
    Method & Clarity & Evidence & Credibility & Actionability & Usefulness \\
    \midrule
    \model & -- & -- & -- & -- & -- \\
    Baseline 1 & -- & -- & -- & -- & -- \\
    Baseline 2 & -- & -- & -- & -- & -- \\
    Baseline 3 & -- & -- & -- & -- & -- \\
    \bottomrule
  \end{tabular}
\end{table}
